% !TeX root = ../main-cn.tex
% UTF-8. Related works synchronized with paper/论文草稿.md.
\section{相关工作}
\label{sec:2_relatedworks}

\subsection{文本到动作生成与预训练基座}

HumanML3D \cite{guo2022humanml3d} 为文本到人体动作生成提供了配对数据与常用评价协议，推动了连续动作建模、潜在空间生成和离散动作 token 建模等路线的发展。
MDM \cite{tevet2023mdm} 与 MotionDiffuse \cite{zhang2024motiondiffuse} 使用扩散过程建模文本条件下的动作分布；MLD \cite{chen2023mld} 则先学习低维动作表示，再在潜在空间中进行扩散，以降低生成的计算开销。
T2M-GPT \cite{zhang2023t2mgpt}、MotionGPT \cite{jiang2023motiongpt} 和 MoMask \cite{guo2024momask} 从离散表示出发，分别探索自回归、语言与动作联合建模以及掩码 token 生成。
LGTM \cite{sun2024lgtm} 研究局部到全局的文本驱动动作建模，进一步增强了细粒度身体语义与整体动作结构之间的联系。

流匹配也为动作生成提供了另一类生成建模范式。
MotionLab \cite{guo2025motionlab} 通过 Motion-Condition-Motion 范式和 rectified flow 统一动作生成与编辑，而 HY-Motion 1.0 \cite{tencent2025hymotion} 探索了大规模流匹配模型在文本动作生成中的扩展能力。
总体而言，现有文本动作生成方法已经形成了扩散、潜在扩散、离散 token 和流匹配等多种成熟基座，为进一步引入外部控制信号提供了较强的动作先验。本文关注如何在尽量保留这些预训练先验的同时，将稀疏 IMU 观测作为额外条件接入生成过程。

\subsection{可控动作生成与条件适配}

可控动作生成在文本语义之外加入轨迹、关键帧或局部关节约束，使用户能够指定动作的空间执行方式。
GMD \cite{karunratanakul2023gmd} 通过特征投影与引导机制结合文本和空间约束，处理轨迹及稀疏关键帧等控制条件。
OmniControl \cite{xie2023omnicontrol} 进一步支持不同时间与不同关节的空间约束，并结合空间引导和真实感引导协调控制精度与动作自然性。
这些工作表明，额外条件能够显著提高动作生成的可控性，同时也需要兼顾条件满足程度与生成动作的自然性。

在预训练模型适配方面，ControlNet \cite{zhang2023controlnet} 在图像扩散模型上展示了通过零初始化控制连接引入额外条件、同时保留原有生成先验的有效性，这一思路也为动作生成中的条件适配提供了参考。
近期 MoCoDiff \cite{song2026mocodiff} 使用模态相关的轻量调制模块处理文本、风格和历史动作条件，进一步关注多条件之间的干扰以及长时序一致性。
因此，将轻量条件模块接入预训练生成模型已经成为可控生成中的常见思路，而不同控制信号的物理含义和信息密度会进一步影响条件表示与融合方式。

与直接给定关节位置或轨迹不同，IMU 通常提供局部加速度和朝向信息，并不显式包含完整的绝对人体轨迹；同时，传感器安装坐标、初始状态和观测部位也会影响可恢复的信息。
这使得稀疏 IMU 控制既包含连续时序约束，又具有明显的局部性和不完备性。
本文据此采用时序表示将 IMU 信号接入文本动作生成过程，并通过不同条件组合考察文本语义与实例级惯性观测之间的互补关系。

\subsection{稀疏惯性人体动作重建}

稀疏惯性动作重建利用少量可穿戴设备估计全身姿态，为低成本人体运动捕捉提供了技术基础。
DIP \cite{huang2018dip} 学习从稀疏惯性观测恢复人体姿态，TransPose \cite{yi2021transpose} 进一步联合估计姿态与人体平移。
PIP \cite{yi2022pip} 在稀疏惯性估计中进一步引入物理约束，将神经运动学估计与物理优化结合，以改善长时序运动的稳定性和物理合理性。
IMUPoser \cite{mollyn2023imuposer} 和 MobilePoser \cite{xu2024mobileposer} 则将这一方向拓展到手机、手表和耳戴设备等消费级硬件，强调稀疏设备条件下的实际可用性。
这些工作为传感器表示、时序建模、全局运动估计以及人体物理约束提供了重要基础。

稀疏惯性重建与本文研究在输入模态上高度相关，但二者的任务目标并不完全相同。
传统重建方法通常以对应真实运动为唯一目标，重点考察姿态、平移和时间稳定性；文本条件动作生成则同时涉及语义匹配、动作分布以及同一语义下的合理多样性。
因此，惯性重建方法能够反映模型从稀疏传感器中恢复人体运动信息的能力，而文本与 IMU 联合生成还需要进一步考察两类条件共同作用时的生成质量和条件一致性。

\subsection{语言与传感器的联合动作建模}

语言与稀疏传感器联合建模与本文最为接近。
Spatial-Related Sensors Matters \cite{yang2024spatialsensors} 利用文本语义辅助惯性动作重建，通过传感器加权、层次时序 Transformer 和对比学习处理传感器特征与语义的对齐。
该工作表明，文本语义能够帮助缓解稀疏惯性观测带来的动作歧义，体现了语言与惯性信号之间的互补性。

Ego4o \cite{wang2025ego4o} 接收不同组合的稀疏 IMU、可选第一视角图像与动作描述，将多模态信息映射到动作 VQ-VAE 的潜在空间，用于动作捕捉与理解。
在不使用视觉输入时，其文本与 IMU 的联合设置与本文具有较直接的输入输出联系。
EgoLM \cite{hong2025egolm} 则借助语言模型联合建模动作与语言，覆盖动作描述以及从文本或稀疏传感器生成动作等任务。
这些研究已经开始将语言、惯性传感器和人体动作放入统一建模框架中，说明多模态条件不仅可以辅助运动重建，也可以用于更一般的动作生成与理解任务。

与面向多模态、多任务的统一系统相比，本文更聚焦于文本与稀疏 IMU 两类条件在预训练文本动作生成基座中的协同作用。
文本主要描述动作类别和语义意图，而 IMU 提供与具体执行实例相关的局部动态观测；二者分别对应较高层的语义约束和较低层的时序运动约束。
因此，如何在保留文本动作先验的同时利用稀疏惯性观测，是本文所研究问题与已有语言辅助惯性重建和多模态动作建模工作的主要联系所在。
